\documentclass[11pt]{article}

\usepackage[margin=1in]{geometry}
\usepackage{amsmath,amssymb,booktabs,tabularx}

\title{Conditional mean baseline clock scale}
\author{}
\date{}

\begin{document}
\maketitle

\noindent\textbf{Status.} Draft\\
\textbf{Plan item.} Conditional mean baseline clock scale \(\overline Z_h\)\\
\textbf{Depends on.} \texttt{01\_long\_variance\_normalization.tex},
\texttt{02\_first\_moment\_guard.tex}\\
\textbf{Feeds.} Long-variance state and return clocks, sufficient-statistics consistency, simulation and aggregation validation

\section{Purpose}

Derive the conditional-mean long-variance scale used by seeded and coarse
simulation when converting deterministic baseline event opportunity into
baseline return-clock exposure.
This explains how simulation uses the same long-variance state as affine
pricing while avoiding pathwise integration of \(Z_t\) inside every coarse
return draw.

\section{Setup}

Consider one interval \(h=[t_h,t_{h+1})\). Let \(\Delta_h^{\mathrm{open}}\) be the elapsed
open-market time over which baseline return-clock exposure is earned.
Conditional on the interval-start state \(Z_{t_h}\), ignore
unobserved trade-event jumps while deriving the mean baseline clock scale.
Those event jumps are applied to the terminal long-variance state through the
event-count update.

Between event jumps, the conditional first moment of \(Z_t\) follows
\[
\frac{d}{ds}\mathbb E[Z_{t_h+s}\mid Z_{t_h}]
=
\kappa_z\left(\bar z-\mathbb E[Z_{t_h+s}\mid Z_{t_h}]\right),
\qquad s\ge 0.
\]

\section{Notation}

\begin{tabularx}{\textwidth}{lXX}
\toprule
Symbol & Meaning & Units or domain\\
\midrule
\(Z_{t_h}\) & Interval-start long-variance state & Positive\\
\(\bar z\) & Derived long-variance mean-reversion level & Positive\\
\(\kappa_z\) & Long-variance mean-reversion rate & Positive inverse time\\
\(\Delta_h^{\mathrm{open}}\) & Elapsed open-market time in interval \(h\) & Non-negative\\
\(\overline Z_h\) & Average conditional-mean long-variance level over open time & Positive\\
\(B_h^\pm\) & Deterministic baseline event opportunity before \(Z\) scaling & Non-negative business time\\
\(B_{Z,h}^\pm\) & Baseline return-clock exposure after \(Z\) scaling & Non-negative business time\\
\bottomrule
\end{tabularx}

\section{Assumptions used}

\begin{tabularx}{\textwidth}{XXX}
\toprule
Assumption & Where it is used & Why it is needed\\
\midrule
The long-variance first-moment dynamics are mean-reverting between jumps.
& Solving the conditional mean \(m(s)\).
& The baseline clock scale is an average conditional first moment.\\
The mean-reversion level \(\bar z\) is positive and derived from unit-mean normalization.
& Keeping \(\overline Z_h\) positive.
& The scale must be a valid positive return-clock multiplier.\\
Baseline return-clock exposure accrues only over open-market elapsed time.
& Integrating over \(\Delta_h^{\mathrm{open}}\).
& Closed time should not create baseline trading-time return exposure.\\
Trade-event jumps are handled by terminal \(Z\) state updates, not by the conditional-mean baseline clock integral.
& Conditioning the clock-scale derivation on the start state and mean-reversion drift.
& This keeps baseline exposure and event-feedback state evolution separate.\\
\bottomrule
\end{tabularx}

\section{Result}

\noindent\textbf{Identity.} The conditional-mean baseline clock scale is
\[
\overline Z_h
=
\begin{cases}
Z_{t_h},
& \Delta_h^{\mathrm{open}}=0,\\[6pt]
\displaystyle
\bar z+
(Z_{t_h}-\bar z)
\frac{1-e^{-\kappa_z\Delta_h^{\mathrm{open}}}}
{\kappa_z\Delta_h^{\mathrm{open}}},
& \Delta_h^{\mathrm{open}}>0.
\end{cases}
\]
Simulation then uses
\[
B_{Z,h}^+
=
\overline Z_h B_h^+,
\qquad
B_{Z,h}^-
=
\overline Z_h B_h^-,
\]
where \(B_h^\pm\) are the deterministic baseline opportunities generated by
seasonality and \(\bar\lambda\) over open time.

\section{Derivation}

Let
\[
m(s)
=
\mathbb E[Z_{t_h+s}\mid Z_{t_h}].
\]
Between event jumps, the conditional first-moment equation is
\[
m'(s)=\kappa_z(\bar z-m(s)),
\qquad
m(0)=Z_{t_h}.
\]
Solving this linear ODE gives
\[
m(s)=\bar z+(Z_{t_h}-\bar z)e^{-\kappa_z s}.
\]

The baseline clock scale is the average conditional mean over open-market
elapsed time:
\[
\overline Z_h
=
\frac{1}{\Delta_h^{\mathrm{open}}}
\int_0^{\Delta_h^{\mathrm{open}}}m(s)\,ds,
\qquad
\Delta_h^{\mathrm{open}}>0.
\]
Substitute the solution for \(m(s)\):
\[
\overline Z_h
=
\frac{1}{\Delta_h^{\mathrm{open}}}
\int_0^{\Delta_h^{\mathrm{open}}}
\left[
\bar z+(Z_{t_h}-\bar z)e^{-\kappa_z s}
\right]ds.
\]
The integral is
\[
\bar z\Delta_h^{\mathrm{open}}
+
(Z_{t_h}-\bar z)
\frac{1-e^{-\kappa_z\Delta_h^{\mathrm{open}}}}{\kappa_z}.
\]
Dividing by \(\Delta_h^{\mathrm{open}}\) gives
\[
\overline Z_h
=
\bar z+
(Z_{t_h}-\bar z)
\frac{1-e^{-\kappa_z\Delta_h^{\mathrm{open}}}}
{\kappa_z\Delta_h^{\mathrm{open}}}.
\]

For \(\Delta_h^{\mathrm{open}}=0\), no open-market baseline return-clock
exposure is earned. The scale is therefore defined as the continuous boundary
value \(Z_{t_h}\), which avoids dividing by zero and leaves any zero-open-time
baseline activity at zero.

\section{Domain and validity}

The identity requires \(Z_{t_h}>0\), \(\bar z>0\), \(\kappa_z>0\), and
\(\Delta_h^{\mathrm{open}}\ge0\).
The formula is a conditional first-moment scale for simulation and coarse
return generation.
It is not a pathwise identity for the stochastic integral
\(\int s(u)\bar\lambda Z_u\,du\).
The affine schedule transform keeps that baseline-clock integral in transform
form through the \(V\)-Riccati equation.

\section{Consequences}

\begin{itemize}
    \item If \(Z_{t_h}=1\) and \(\bar z=1\), then \(\overline Z_h=1\) for every open-time length.
    \item If \(Z_{t_h}>\bar z\), the baseline clock scale is above \(\bar z\) but decays toward it as open time grows.
    \item If \(Z_{t_h}<\bar z\), the baseline clock scale is below \(\bar z\) but increases toward it as open time grows.
    \item Simulation can use \(\overline Z_h\) for baseline return-clock exposure while separately sampling terminal \(Z\) and applying event-count jumps.
\end{itemize}

\section{Checks}

\begin{tabularx}{\textwidth}{XXX}
\toprule
Check & Expected result & Evidence location\\
\midrule
Kernel baseline clock scale
& The implementation returns \(Z_{t_h}\) at zero elapsed time and the integrated conditional mean divided by elapsed time otherwise.
& \texttt{LongVarianceState::baseline\_clock\_scale} in \texttt{crates/janiform-model-kernels/src/long\_variance.rs}\\
Neutral transition
& Unit long variance with neutral parameters keeps baseline clock scale equal to one.
& Test \texttt{neutral\_long\_variance\_transition\_keeps\_unit\_clock\_and\_state}\\
Elevated initial state
& An elevated initial \(Z\) produces a baseline clock scale above one under mean reversion to one.
& Test \texttt{elevated\_initial\_state\_increases\_baseline\_clock\_scale}\\
Simulation routing
& Coarse trading-time return generation receives the sampled transition's baseline long-variance clock scale.
& \texttt{sample\_transition\_trading\_time\_return} call in \texttt{crates/janiform-engine-simulation/src/runtime/simulation.rs}\\
\bottomrule
\end{tabularx}

\end{document}
